%% French through polyglossia rather than babel: the judgment half of
%% babel-fr.tex, and the other side of its \fg half.
%%
%% polyglossia leaves ? ! : ; alone as characters and adds the thin space
%% before them by other means -- interchar tokens under XeTeX, a node
%% callback under LuaTeX -- so the judgment scanner meets an ordinary ?
%% here, and French spacing is inserted around marks it has collected.  A
%% mark must still hang and be collected whole.
%%
%% polyglossia defines no \og/\fg, so \fg is free and linguexx claims \fg.
%% for its sixth glossed sub-example, as it does in an English document.
%% babel-fr asserts that babel's \fg survives; this asserts the shorthand
%% is there when nothing owns the name.  polyglossia is loaded AFTER
%% linguexx; rtl-polyglossia loads it before.
\documentclass[11pt]{article}
\usepackage[a4paper,margin=25mm]{geometry}
\usepackage{iftex}
\ifPDFTeX
  \usepackage[T1]{fontenc}
  \usepackage[utf8]{inputenc}
\else
  \usepackage{fontspec}
\fi
\usepackage{linguexx}
\pagestyle{empty}

\usepackage{polyglossia}
\setdefaultlanguage{french}
\begin{document}
\ex. FRMAIN une phrase correcte.
\a. FRPLAIN sous-exemple correct.
\b. ?FRQ un point d'interrogation.
\c. ??FRQQ deux points, collectes ensemble.
\d. ?*FRQS un point et une etoile.
\e. *FRST une etoile.
\z.

\ex. FRFINAL Est-ce que cette phrase est correcte ?

\ex. FRGLOSS
\ag. AGA le chat \\ AGB the cat \\
\bg. BGA le chien \\ BGB the dog \\
\cg. CGA la vache \\ CGB the cow \\
\dg. DGA le loup \\ DGB the wolf \\
\eg. EGA la chouette \\ EGB the owl \\
\fg. FGA la souris \\ FGB the mouse \\
\z.
\end{document}
